\documentclass{article}
\usepackage[margin=25mm]{geometry}
\usepackage{amsmath,amssymb,xurl}
\usepackage[hidelinks]{hyperref}
% Shared commands for the notes in docs/paper-gaps/.
% Every note loads this file via % Shared commands for the notes in docs/paper-gaps/.
% Every note loads this file via % Shared commands for the notes in docs/paper-gaps/.
% Every note loads this file via \input{command}. Project-wide notation
% belongs here, never in per-note redefinitions.

% --- Project configuration ------------------------------------------------
\newcommand{\projectIssueBase}{https://github.com/LionSR/QICLean/issues/}
\newcommand{\projectPrBase}{https://github.com/LionSR/QICLean/pull/}
\newcommand{\projectDocBase}{https://sirui-lu.com/QICLean/docs/}
% ---------------------------------------------------------------------------

\DeclareUnicodeCharacter{2080}{\ifmmode{}_{0}\else\textsubscript{0}\fi}
\DeclareUnicodeCharacter{2081}{\ifmmode{}_{1}\else\textsubscript{1}\fi}
\DeclareUnicodeCharacter{2082}{\ifmmode{}_{2}\else\textsubscript{2}\fi}
\DeclareUnicodeCharacter{2099}{\ifmmode{}_{n}\else\textsubscript{n}\fi}

\newcommand{\Lean}{\textsc{Lean}}

% A formal declaration name, upright sans, breakable at underscores.
\ExplSyntaxOn
\NewDocumentCommand{\leanid}{m}{
  \ifmmode
    \textsf{\detokenize{#1}}
  \else
    \group_begin:
    \urlstyle{sf}
    \tl_set:Nn \l_tmpa_tl {#1}
    \tl_replace_all:Nnn \l_tmpa_tl {\_} {_}
    \exp_args:NV \path \l_tmpa_tl
    \group_end:
  \fi
}
\ExplSyntaxOff

% Traceability links; use in footnotes, not in the running prose.
\newcommand{\ghissue}[1]{\href{\projectIssueBase#1}{GitHub issue \##1}}
\newcommand{\ghpr}[1]{\href{\projectPrBase#1}{PR \##1}}
\newcommand{\doclink}[1]{\href{\projectDocBase#1.html}{\path{#1}}}

% Domain notation, using \providecommand so a note-local refinement stays
% possible.
\providecommand{\tr}{\operatorname{tr}}

% Dirac notation and the identity operator, as in blueprint/src/macros.
\usepackage{dsfont}
\providecommand{\ket}[1]{|#1\rangle}
\providecommand{\bra}[1]{\langle #1|}
\providecommand{\braket}[2]{\langle #1 | #2 \rangle}
\providecommand{\ketbra}[2]{|#1\rangle\!\langle #2|}
\providecommand{\Id}{\mathds{1}}
\providecommand{\one}{\mathds{1}}
\providecommand{\C}{\mathbb{C}}
\providecommand{\MN}[1]{M_{#1}(\C)}
\providecommand{\MD}{M_D(\C)}

% Verdict marker: \gapnote{<kind>}{<status>}. Kind is the policy
% classification (clarification, local-correction, scope-restriction,
% unfaithful, false-source, open-gap); status is open, wip (elimination
% underway), resolved, or historical. Machine-read by the site index;
% prints a verdict line.
\newcommand{\gapnote}[2]{\par\noindent\textbf{Verdict:} #1 (#2).\par}
. Project-wide notation
% belongs here, never in per-note redefinitions.

% --- Project configuration ------------------------------------------------
\newcommand{\projectIssueBase}{https://github.com/LionSR/QICLean/issues/}
\newcommand{\projectPrBase}{https://github.com/LionSR/QICLean/pull/}
\newcommand{\projectDocBase}{https://sirui-lu.com/QICLean/docs/}
% ---------------------------------------------------------------------------

\DeclareUnicodeCharacter{2080}{\ifmmode{}_{0}\else\textsubscript{0}\fi}
\DeclareUnicodeCharacter{2081}{\ifmmode{}_{1}\else\textsubscript{1}\fi}
\DeclareUnicodeCharacter{2082}{\ifmmode{}_{2}\else\textsubscript{2}\fi}
\DeclareUnicodeCharacter{2099}{\ifmmode{}_{n}\else\textsubscript{n}\fi}

\newcommand{\Lean}{\textsc{Lean}}

% A formal declaration name, upright sans, breakable at underscores.
\ExplSyntaxOn
\NewDocumentCommand{\leanid}{m}{
  \ifmmode
    \textsf{\detokenize{#1}}
  \else
    \group_begin:
    \urlstyle{sf}
    \tl_set:Nn \l_tmpa_tl {#1}
    \tl_replace_all:Nnn \l_tmpa_tl {\_} {_}
    \exp_args:NV \path \l_tmpa_tl
    \group_end:
  \fi
}
\ExplSyntaxOff

% Traceability links; use in footnotes, not in the running prose.
\newcommand{\ghissue}[1]{\href{\projectIssueBase#1}{GitHub issue \##1}}
\newcommand{\ghpr}[1]{\href{\projectPrBase#1}{PR \##1}}
\newcommand{\doclink}[1]{\href{\projectDocBase#1.html}{\path{#1}}}

% Domain notation, using \providecommand so a note-local refinement stays
% possible.
\providecommand{\tr}{\operatorname{tr}}

% Dirac notation and the identity operator, as in blueprint/src/macros.
\usepackage{dsfont}
\providecommand{\ket}[1]{|#1\rangle}
\providecommand{\bra}[1]{\langle #1|}
\providecommand{\braket}[2]{\langle #1 | #2 \rangle}
\providecommand{\ketbra}[2]{|#1\rangle\!\langle #2|}
\providecommand{\Id}{\mathds{1}}
\providecommand{\one}{\mathds{1}}
\providecommand{\C}{\mathbb{C}}
\providecommand{\MN}[1]{M_{#1}(\C)}
\providecommand{\MD}{M_D(\C)}

% Verdict marker: \gapnote{<kind>}{<status>}. Kind is the policy
% classification (clarification, local-correction, scope-restriction,
% unfaithful, false-source, open-gap); status is open, wip (elimination
% underway), resolved, or historical. Machine-read by the site index;
% prints a verdict line.
\newcommand{\gapnote}[2]{\par\noindent\textbf{Verdict:} #1 (#2).\par}
. Project-wide notation
% belongs here, never in per-note redefinitions.

% --- Project configuration ------------------------------------------------
\newcommand{\projectIssueBase}{https://github.com/LionSR/QICLean/issues/}
\newcommand{\projectPrBase}{https://github.com/LionSR/QICLean/pull/}
\newcommand{\projectDocBase}{https://sirui-lu.com/QICLean/docs/}
% ---------------------------------------------------------------------------

\DeclareUnicodeCharacter{2080}{\ifmmode{}_{0}\else\textsubscript{0}\fi}
\DeclareUnicodeCharacter{2081}{\ifmmode{}_{1}\else\textsubscript{1}\fi}
\DeclareUnicodeCharacter{2082}{\ifmmode{}_{2}\else\textsubscript{2}\fi}
\DeclareUnicodeCharacter{2099}{\ifmmode{}_{n}\else\textsubscript{n}\fi}

\newcommand{\Lean}{\textsc{Lean}}

% A formal declaration name, upright sans, breakable at underscores.
\ExplSyntaxOn
\NewDocumentCommand{\leanid}{m}{
  \ifmmode
    \textsf{\detokenize{#1}}
  \else
    \group_begin:
    \urlstyle{sf}
    \tl_set:Nn \l_tmpa_tl {#1}
    \tl_replace_all:Nnn \l_tmpa_tl {\_} {_}
    \exp_args:NV \path \l_tmpa_tl
    \group_end:
  \fi
}
\ExplSyntaxOff

% Traceability links; use in footnotes, not in the running prose.
\newcommand{\ghissue}[1]{\href{\projectIssueBase#1}{GitHub issue \##1}}
\newcommand{\ghpr}[1]{\href{\projectPrBase#1}{PR \##1}}
\newcommand{\doclink}[1]{\href{\projectDocBase#1.html}{\path{#1}}}

% Domain notation, using \providecommand so a note-local refinement stays
% possible.
\providecommand{\tr}{\operatorname{tr}}

% Dirac notation and the identity operator, as in blueprint/src/macros.
\usepackage{dsfont}
\providecommand{\ket}[1]{|#1\rangle}
\providecommand{\bra}[1]{\langle #1|}
\providecommand{\braket}[2]{\langle #1 | #2 \rangle}
\providecommand{\ketbra}[2]{|#1\rangle\!\langle #2|}
\providecommand{\Id}{\mathds{1}}
\providecommand{\one}{\mathds{1}}
\providecommand{\C}{\mathbb{C}}
\providecommand{\MN}[1]{M_{#1}(\C)}
\providecommand{\MD}{M_D(\C)}

% Verdict marker: \gapnote{<kind>}{<status>}. Kind is the policy
% classification (clarification, local-correction, scope-restriction,
% unfaithful, false-source, open-gap); status is open, wip (elimination
% underway), resolved, or historical. Machine-read by the site index;
% prints a verdict line.
\newcommand{\gapnote}[2]{\par\noindent\textbf{Verdict:} #1 (#2).\par}

\title{Strict dimension decrease in the low-density bootstrap}
\author{}
\date{2026-11-02}
\begin{document}
\maketitle
\gapnote{local-correction}{resolved}

\section{The minimal dimension is a natural number}

Arad, Landau, Vazirani and Vidick,
\emph{Rigorous RG algorithms and area laws for low energy eigenstates in 1D},
\href{https://arxiv.org/abs/1602.08828}{arXiv:1602.08828v2},
Proposition \texttt{prop:albootstrap-ld}, source lines 1112--1128,
uses the minimal dimension of a $0.015$-viable subspace for each nested
low-energy target. These are minima over natural dimensions, with the
whole regional space supplying a competitor.

Write $p$ for the minimum at the lower target and $Q$ for the minimum
at the next upper target. The proof at lines 1115 and 1125 uses a sample
of size $p/D^2$; applying $D^2$ operators then bounds the image dimension
by $p$. This bound alone does not contradict minimality.
Use instead
\[
    s=\left\lfloor\frac{p}{2D^2}\right\rfloor.
\]
The actual image of an $s$-dimensional sample under $D^2$ operators has
\[
    \dim\operatorname{span}_i A_iS\le D^2s\le\lfloor p/2\rfloor<p
\]
when $p>0$. In the nontrivial sampling branch $p\ge4D^2$, the quotient
is nonzero and satisfies $s\ge p/(4D^2)$. The small-dimension branch
must be settled before sampling.

\section{Transfer between the two targets}

Here $U$ is $\varepsilon$-close to $T$ when every $t\in T$ satisfies
$\|P_Ut\|^2\ge(1-\varepsilon)\|t\|^2$. This fixes the direction of the
following two transfers.

Let $G_+,G_-$ be the original upper and lower targets, and let
$T_{\rm hi},T_{\rm lo}$ be the targets for the improving family.
Suppose $G_+$ is $1/200$-close to $T_{\rm hi}$ and $T_{\rm lo}$ is
$1/200$-close to $G_-$. A $3/200$-viable upper input is then
$1/25$-viable for $T_{\rm hi}$. A successful sample has overlap parameter
\[
    \alpha=\frac{(1-1/25)s}{8Q}=\frac{3s}{25Q}.
\]
If the same family improves its sample error to $\Delta/\alpha^2$,
transfer to the lower original target gives error
\[
    2\left(\frac{\Delta}{\alpha^2}+\frac1{200}\right).
\]
Thus $\Delta/\alpha^2\le1/400$ gives the required $3/200$-viability.
The sample is chosen toward $T_{\rm hi}$, where the improving-family
hypothesis applies, before the final transfer to $G_-$.

\section{One finite numerical calibration}

A sufficient scalar calibration is
\[
 \log D\ge100,\quad r\ge1,\quad
 \log D'\le\frac{13}{6}\log D,\quad
 rD^3<p\le Q\le r(D')^3,\quad D^{12}\Delta\le10^{-5}.
\]
It implies $D\ge8$ and $p\ge4D^2$. With the actual natural quotient above,
\[
 \left(1+\frac2{\sqrt\alpha}\right)^{2t}Qe^{-s/16}<1,
 \qquad \frac{\Delta}{\alpha^2}\le\frac1{400}
 \quad\text{for every integer }t\le r.
\]
The exponent $2t$ is the corrected complex-sphere exponent, recorded in
\path{alvv16_random_sampling_net_count.tex}.
The count comparison gives $Q^2\le p^2D^7$, hence
$\alpha^2\ge9/(10000D^{11})$. Also $s>rD/4$ and
$1+2/\sqrt\alpha\le13D^3$. Writing $x=\log D\ge100$, the logarithm
of the failure expression is at most $13rx-rD/64<0$, since
$D=e^x\ge x^3/6>832x$. The resulting sampling step between two nested
targets is \leanid{Submodule.exists_isViable_nested_sample_of_count_comparison},
and the induction over nested improving families that selects the natural
minima is \leanid{Submodule.exists_isViable_finrank_le_of_dyadic_families}.

The count and suppression conditions must be proved together for the same
physical family. Increasing a count bound after proving its suppression
tradeoff does not establish these conditions. The physical spectral-window
comparisons and the complete low-density induction are separate from this
finite sampling repair.

\section{A consistent positive logarithm}

Source line 1112 uses a power $\log^c(\log n/\mu)$, whereas the
comparison function in line 1113 displays only one logarithm. The same
power must be retained throughout the comparison. At finite scales,
use a positive regularized expression instead of a logarithm that can
vanish or be negative. For a fixed bandwidth $\eta>0$ and deficit
$v=\mu/\log n$, the actual count admits the explicit bound
\begin{align*}
 \log(D_0^3)\le \frac{K_\eta}{v}
     \bigl(1+\log d+\log(1+1/v)\bigr)^3.
\end{align*}
The coefficient is chosen before the physical length, local dimension
and deficit. This records the finite logarithmic convention used with
the source's soft-O notation, not a claim that its area-law theorem fails.
\end{document}
